\chapter{Código Equivalente por Plataforma}
\label{apendice:equivalencias}
% ── índice ──────────────────────────────────────────────────────
\index{equivalências Stata/R/Python|textbf}

Este apêndice reproduz DGPs representativos em Stata, R e Python, lado a
lado com o código \hayashi{}. O objetivo é duplo: permitir validação cruzada
de resultados e evidenciar diferenças de sintaxe.

\begin{nota}
  Os geradores de números aleatórios diferem entre plataformas — os
  valores numéricos serão distintos, mas as propriedades estatísticas
  (recuperação de parâmetros, direção do viés, decisão do Hausman)
  são idênticas para qualquer semente.
\end{nota}

% ─────────────────────────────────────────────────────────────────────────────
\section{OLS — DGP sem ruído e comparação com variável omitida}

\subsection*{Hayashi}

\begin{lstlisting}
seed(42)
input df
id
1
end
for i in 1..=10 { df = append(df, df) }
df |> filter(_n <= 1000)
generate df x1 = rnormal()
generate df x2 = rnormal()
generate df y  = 2.5 + 1.8*x1 - 0.7*x2
let m1 = ols(y ~ x1 + x2, df)
let m2 = ols(y ~ x1, df)
esttab(m1, m2)
\end{lstlisting}

\subsection*{Stata}

\begin{lstlisting}[language={}, basicstyle=\ttfamily\small,
  frame=single, numbers=none, backgroundcolor=\color{codbg}]
clear
set obs 1000
set seed 42
gen x1 = rnormal()
gen x2 = rnormal()
gen y  = 2.5 + 1.8*x1 - 0.7*x2
reg y x1 x2
estimates store m1
reg y x1
estimates store m2
esttab m1 m2          // ssc install estout
\end{lstlisting}

\subsection*{R}

\begin{lstlisting}[language=R, basicstyle=\ttfamily\small,
  frame=single, numbers=none, backgroundcolor=\color{codbg},
  keywordstyle=\color{keyword}\bfseries,
  stringstyle=\color{string},
  commentstyle=\color{comment}\itshape]
set.seed(42)
n  <- 1000
x1 <- rnorm(n);  x2 <- rnorm(n)
y  <- 2.5 + 1.8*x1 - 0.7*x2
m1 <- lm(y ~ x1 + x2)
m2 <- lm(y ~ x1)
library(stargazer)
stargazer(m1, m2, type = "text")
\end{lstlisting}

\subsection*{Python}

\begin{lstlisting}[language=Python, basicstyle=\ttfamily\small,
  frame=single, numbers=none, backgroundcolor=\color{codbg},
  keywordstyle=\color{keyword}\bfseries,
  stringstyle=\color{string},
  commentstyle=\color{comment}\itshape]
import numpy as np, pandas as pd
import statsmodels.formula.api as smf
from statsmodels.iolib.summary2 import summary_col
np.random.seed(42)
n = 1000
x1 = np.random.randn(n);  x2 = np.random.randn(n)
y  = 2.5 + 1.8*x1 - 0.7*x2
df = pd.DataFrame({'y': y, 'x1': x1, 'x2': x2})
m1 = smf.ols('y ~ x1 + x2', data=df).fit()
m2 = smf.ols('y ~ x1',      data=df).fit()
print(summary_col([m1, m2], stars=True))
\end{lstlisting}

% ─────────────────────────────────────────────────────────────────────────────
\section{IV — Endogeneidade com instrumento válido}

\subsection*{Hayashi}

\begin{lstlisting}
seed(42)
...
generate df z   = rnormal()
generate df v   = rnormal()
let d1          = cov(df, z, v) / variance(df, z)
let d0          = mean(df, v) - d1 * mean(df, z)
generate df eps = v - d0 - d1*z
generate df x   = 0.8*z + eps
generate df y   = 1.5 + 2.0*x + 0.7*eps
let m_ols = ols(y ~ x, df)
let m_iv  = iv(y ~ x, ~ z, df)
esttab(m_ols, m_iv)
\end{lstlisting}

\subsection*{Stata}

\begin{lstlisting}[language={}, basicstyle=\ttfamily\small,
  frame=single, numbers=none, backgroundcolor=\color{codbg}]
gen z = rnormal()
gen v = rnormal()
reg v z
predict eps, residuals
gen x = 0.8*z + eps
gen y = 1.5 + 2.0*x + 0.7*eps
reg y x
ivregress 2sls y (x = z)
\end{lstlisting}

\subsection*{R}

\begin{lstlisting}[language=R, basicstyle=\ttfamily\small,
  frame=single, numbers=none, backgroundcolor=\color{codbg},
  keywordstyle=\color{keyword}\bfseries,
  stringstyle=\color{string},
  commentstyle=\color{comment}\itshape]
eps <- resid(lm(v ~ z))
x   <- 0.8*z + eps
y   <- 1.5 + 2.0*x + 0.7*eps
library(ivreg)
ivreg(y ~ x | z, data = data.frame(y,x,z))
\end{lstlisting}

\subsection*{Python}

\begin{lstlisting}[language=Python, basicstyle=\ttfamily\small,
  frame=single, numbers=none, backgroundcolor=\color{codbg},
  keywordstyle=\color{keyword}\bfseries,
  stringstyle=\color{string},
  commentstyle=\color{comment}\itshape]
from linearmodels.iv import IV2SLS
import statsmodels.api as sm
IV2SLS(y, sm.add_constant(x),
       None, sm.add_constant(z)).fit()
\end{lstlisting}

% ─────────────────────────────────────────────────────────────────────────────
\section{Painel — FE, RE e Hausman}

\subsection*{Hayashi}

\begin{lstlisting}
seed(42) / ... / df |> filter(_n <= 1000)
generate df id    = (_n - 1) % 100 + 1
generate df alpha = id * 0.05
generate df x     = alpha + rnormal()
generate df y     = 1.0 + 2.0*x + alpha + rnormal()
let m_fe  = fe(y ~ x, df, id=id)
let m_re  = re(y ~ x, df, id=id)
hausman(m_fe, m_re)
\end{lstlisting}

\subsection*{Stata}

\begin{lstlisting}[language={}, basicstyle=\ttfamily\small,
  frame=single, numbers=none, backgroundcolor=\color{codbg}]
xtset id
xtreg y x, fe
estimates store m_fe
xtreg y x, re
estimates store m_re
hausman m_fe m_re
\end{lstlisting}

\subsection*{R}

\begin{lstlisting}[language=R, basicstyle=\ttfamily\small,
  frame=single, numbers=none, backgroundcolor=\color{codbg},
  keywordstyle=\color{keyword}\bfseries,
  stringstyle=\color{string},
  commentstyle=\color{comment}\itshape]
library(plm)
pd   <- pdata.frame(df, index = "id")
m_fe <- plm(y ~ x, data = pd, model = "within")
m_re <- plm(y ~ x, data = pd, model = "random")
phtest(m_fe, m_re)
\end{lstlisting}

\subsection*{Python}

\begin{lstlisting}[language=Python, basicstyle=\ttfamily\small,
  frame=single, numbers=none, backgroundcolor=\color{codbg},
  keywordstyle=\color{keyword}\bfseries,
  stringstyle=\color{string},
  commentstyle=\color{comment}\itshape]
from linearmodels import PanelOLS, RandomEffects
m_fe = PanelOLS(df['y'], df[['x']], entity_effects=True).fit()
m_re = RandomEffects(df['y'], df[['x']]).fit()
\end{lstlisting}

% ─────────────────────────────────────────────────────────────────────────────
\section{Logit e Probit}

\subsection*{Hayashi}

\begin{lstlisting}
generate df p = 1 / (1 + exp(-(-1.0 + 2.0*x)))
generate df y = runiform() < p
let m_l = logit(y ~ x, df)
let m_p = probit(y ~ x, df)
margins(m_l, df)
\end{lstlisting}

\subsection*{Stata}

\begin{lstlisting}[language={}, basicstyle=\ttfamily\small,
  frame=single, numbers=none, backgroundcolor=\color{codbg}]
logit  y x
probit y x
margins, dydx(x)
\end{lstlisting}

\subsection*{R}

\begin{lstlisting}[language=R, basicstyle=\ttfamily\small,
  frame=single, numbers=none, backgroundcolor=\color{codbg},
  keywordstyle=\color{keyword}\bfseries,
  stringstyle=\color{string},
  commentstyle=\color{comment}\itshape]
m_l <- glm(y ~ x, family = binomial("logit"))
m_p <- glm(y ~ x, family = binomial("probit"))
library(margins)
summary(margins(m_l))
\end{lstlisting}

\subsection*{Python}

\begin{lstlisting}[language=Python, basicstyle=\ttfamily\small,
  frame=single, numbers=none, backgroundcolor=\color{codbg},
  keywordstyle=\color{keyword}\bfseries,
  stringstyle=\color{string},
  commentstyle=\color{comment}\itshape]
import statsmodels.formula.api as smf
m_l = smf.logit('y ~ x', data=df).fit()
print(m_l.get_margeff().summary())
\end{lstlisting}

% ─────────────────────────────────────────────────────────────────────────────
\section{Modelos de contagem}

\subsection*{Hayashi}

\begin{lstlisting}
generate df lam = exp(0.5 + 1.0*x)
generate df y   = floor(lam + runiform())
let m_p  = poisson(y ~ x, df)
let m_nb = nbreg(y ~ x, df)
\end{lstlisting}

\subsection*{Stata}

\begin{lstlisting}[language={}, basicstyle=\ttfamily\small,
  frame=single, numbers=none, backgroundcolor=\color{codbg}]
poisson y x
nbreg   y x
\end{lstlisting}

\subsection*{R}

\begin{lstlisting}[language=R, basicstyle=\ttfamily\small,
  frame=single, numbers=none, backgroundcolor=\color{codbg},
  keywordstyle=\color{keyword}\bfseries,
  stringstyle=\color{string},
  commentstyle=\color{comment}\itshape]
glm(y ~ x, family = poisson("log"))
library(MASS)
glm.nb(y ~ x)
\end{lstlisting}

\subsection*{Python}

\begin{lstlisting}[language=Python, basicstyle=\ttfamily\small,
  frame=single, numbers=none, backgroundcolor=\color{codbg},
  keywordstyle=\color{keyword}\bfseries,
  stringstyle=\color{string},
  commentstyle=\color{comment}\itshape]
smf.poisson('y ~ x', data=df).fit()
smf.negativebinomial('y ~ x', data=df).fit()
\end{lstlisting}

% ─────────────────────────────────────────────────────────────────────────────
\section{Tobit e Heckman}

\subsection*{Hayashi}

\begin{lstlisting}
let m_t = tobit(y_tobit ~ x, df)
let m_h = heckman(y_out ~ x, s ~ z, df)
\end{lstlisting}

\subsection*{Stata}

\begin{lstlisting}[language={}, basicstyle=\ttfamily\small,
  frame=single, numbers=none, backgroundcolor=\color{codbg}]
tobit y_tobit x, ll(0)
heckman y_out x, select(s = x z)
\end{lstlisting}

\subsection*{R}

\begin{lstlisting}[language=R, basicstyle=\ttfamily\small,
  frame=single, numbers=none, backgroundcolor=\color{codbg},
  keywordstyle=\color{keyword}\bfseries,
  stringstyle=\color{string},
  commentstyle=\color{comment}\itshape]
library(AER)
tobit(y_tobit ~ x, left = 0)
library(sampleSelection)
heckit(s ~ x + z, y_out ~ x)
\end{lstlisting}

\subsection*{Python}

\begin{lstlisting}[language=Python, basicstyle=\ttfamily\small,
  frame=single, numbers=none, backgroundcolor=\color{codbg},
  keywordstyle=\color{keyword}\bfseries,
  stringstyle=\color{string},
  commentstyle=\color{comment}\itshape]
# Tobit: sem implementação nativa no statsmodels
# Heckman: via heckit do pyheckit ou statsmodels
from statsmodels.regression.linear_model import OLS
# implementação manual via two-step
\end{lstlisting}

% ─────────────────────────────────────────────────────────────────────────────
\section{Modelos multinomiais ordenados}

\subsection*{Hayashi}

\begin{lstlisting}
let m_ol = ologit(y ~ x, df)
let m_op = oprobit(y ~ x, df)
let m_ml = mlogit(y ~ x, df)
\end{lstlisting}

\subsection*{Stata}

\begin{lstlisting}[language={}, basicstyle=\ttfamily\small,
  frame=single, numbers=none, backgroundcolor=\color{codbg}]
ologit  y x
oprobit y x
mlogit  y x
\end{lstlisting}

\subsection*{R}

\begin{lstlisting}[language=R, basicstyle=\ttfamily\small,
  frame=single, numbers=none, backgroundcolor=\color{codbg},
  keywordstyle=\color{keyword}\bfseries,
  stringstyle=\color{string},
  commentstyle=\color{comment}\itshape]
library(MASS)
polr(factor(y) ~ x, method = "logistic")  # ologit
polr(factor(y) ~ x, method = "probit")    # oprobit
library(nnet)
multinom(y ~ x)                           # mlogit
\end{lstlisting}

\subsection*{Python}

\begin{lstlisting}[language=Python, basicstyle=\ttfamily\small,
  frame=single, numbers=none, backgroundcolor=\color{codbg},
  keywordstyle=\color{keyword}\bfseries,
  stringstyle=\color{string},
  commentstyle=\color{comment}\itshape]
from statsmodels.miscmodels.ordinal_model import OrderedModel
OrderedModel(y, X, distr='logit').fit()
OrderedModel(y, X, distr='probit').fit()
smf.mnlogit('y ~ x', data=df).fit()
\end{lstlisting}

% ─────────────────────────────────────────────────────────────────────────────
\section{Regularização (LASSO, Ridge, Elastic Net)}

\subsection*{Hayashi}

\begin{lstlisting}
let m_las = lasso(y ~ x1+x2+...+x10, df, alpha=0.1)
let m_rid = ridge_reg(y ~ x1+...+x10, df, alpha=0.5)
let m_ent = enet(y ~ x1+...+x10, df, alpha=0.5)
display m_las
\end{lstlisting}

\subsection*{Stata}

\begin{lstlisting}[language={}, basicstyle=\ttfamily\small,
  frame=single, numbers=none, backgroundcolor=\color{codbg}]
// requer Stata 16+
lasso linear y x1-x10
cvlasso y x1-x10
\end{lstlisting}

\subsection*{R}

\begin{lstlisting}[language=R, basicstyle=\ttfamily\small,
  frame=single, numbers=none, backgroundcolor=\color{codbg},
  keywordstyle=\color{keyword}\bfseries,
  stringstyle=\color{string},
  commentstyle=\color{comment}\itshape]
library(glmnet)
X <- model.matrix(y ~ ., df)[,-1]
# LASSO (alpha=1), Ridge (alpha=0), Enet (alpha=0.5)
cv.glmnet(X, df$y, alpha = 1)   # LASSO com CV
cv.glmnet(X, df$y, alpha = 0)   # Ridge com CV
cv.glmnet(X, df$y, alpha = 0.5) # Enet com CV
\end{lstlisting}

\subsection*{Python}

\begin{lstlisting}[language=Python, basicstyle=\ttfamily\small,
  frame=single, numbers=none, backgroundcolor=\color{codbg},
  keywordstyle=\color{keyword}\bfseries,
  stringstyle=\color{string},
  commentstyle=\color{comment}\itshape]
from sklearn.linear_model import Lasso, Ridge, ElasticNet
Lasso(alpha=0.1).fit(X, y)
Ridge(alpha=0.5).fit(X, y)
ElasticNet(alpha=0.1, l1_ratio=0.5).fit(X, y)
\end{lstlisting}

{\footnotesize \textbf{Nota de escala:} \texttt{alpha} em Hayashi e statsmodels é $\lambda$
(força de penalização). Em glmnet (R), os dados são padronizados automaticamente
e \texttt{lambda} é escolhido por validação cruzada. Em scikit-learn, \texttt{alpha}
é análogo a $\lambda$.}

% ─────────────────────────────────────────────────────────────────────────────
\section{Análise de sobrevivência}

\subsection*{Hayashi}

\begin{lstlisting}
let m_km  = km(time, event, df)
let m_cox = cox(time ~ x, df, event=event)
display m_cox
\end{lstlisting}

\subsection*{Stata}

\begin{lstlisting}[language={}, basicstyle=\ttfamily\small,
  frame=single, numbers=none, backgroundcolor=\color{codbg}]
stset time, failure(event)
sts graph              // Kaplan-Meier
stcox x                // Cox PH
stcox x, nohr          // coeficientes log-hazard
\end{lstlisting}

\subsection*{R}

\begin{lstlisting}[language=R, basicstyle=\ttfamily\small,
  frame=single, numbers=none, backgroundcolor=\color{codbg},
  keywordstyle=\color{keyword}\bfseries,
  stringstyle=\color{string},
  commentstyle=\color{comment}\itshape]
library(survival)
km_fit  <- survfit(Surv(time, event) ~ 1)
cox_fit <- coxph(Surv(time, event) ~ x)
summary(cox_fit)
\end{lstlisting}

\subsection*{Python}

\begin{lstlisting}[language=Python, basicstyle=\ttfamily\small,
  frame=single, numbers=none, backgroundcolor=\color{codbg},
  keywordstyle=\color{keyword}\bfseries,
  stringstyle=\color{string},
  commentstyle=\color{comment}\itshape]
from lifelines import KaplanMeierFitter, CoxPHFitter
kmf = KaplanMeierFitter().fit(df['time'], df['event'])
cph = CoxPHFitter().fit(df, 'time', 'event')
\end{lstlisting}

% ─────────────────────────────────────────────────────────────────────────────
\section{SVAR — Cholesky e Blanchard-Quah}

\subsection*{Hayashi}

\begin{lstlisting}
let m_ch = svar(df, y1, y2, lags=1, id=cholesky)
let m_bq = svar(df, y1, y2, lags=1, id=longrun)
let m_ir = sirf(m_ch, steps=12)
display m_ir
\end{lstlisting}

\subsection*{Stata}

\begin{lstlisting}[language={}, basicstyle=\ttfamily\small,
  frame=single, numbers=none, backgroundcolor=\color{codbg}]
var y1 y2, lags(1)
// Cholesky (padrao do IRF pos-VAR):
irf create myirf, step(12) replace
irf table sirf
// Blanchard-Quah requer restrições manuais via svar
matrix A = (1, 0 \ ., 1)
matrix B = (., 0 \ ., .)
svar y1 y2, lags(1) aeq(A) beq(B)
\end{lstlisting}

\subsection*{R}

\begin{lstlisting}[language=R, basicstyle=\ttfamily\small,
  frame=single, numbers=none, backgroundcolor=\color{codbg},
  keywordstyle=\color{keyword}\bfseries,
  stringstyle=\color{string},
  commentstyle=\color{comment}\itshape]
library(vars)
m <- VAR(cbind(y1, y2), p = 1)
# Cholesky:
irf(m, n.ahead = 12, ortho = TRUE)
# Blanchard-Quah:
BQ(m)
\end{lstlisting}

\subsection*{Python}

\begin{lstlisting}[language=Python, basicstyle=\ttfamily\small,
  frame=single, numbers=none, backgroundcolor=\color{codbg},
  keywordstyle=\color{keyword}\bfseries,
  stringstyle=\color{string},
  commentstyle=\color{comment}\itshape]
from statsmodels.tsa.api import VAR
m  = VAR(df[['y1', 'y2']]).fit(1)
# Cholesky:
m.irf(12).orth_ma            # orthogonalized IRF
# BQ: requer implementação manual em Python
\end{lstlisting}

% ─────────────────────────────────────────────────────────────────────────────
\section{Filtros de decomposição}

\subsection*{Hayashi}

\begin{lstlisting}
hprescott(df, y, lambda=1600)
baxter_king(df, y, low=6, high=32, k=12)
stl(df, y, period=12)
christiano_fitzgerald(df, y, low=6, high=32)
\end{lstlisting}

\subsection*{Stata}

\begin{lstlisting}[language={}, basicstyle=\ttfamily\small,
  frame=single, numbers=none, backgroundcolor=\color{codbg}]
// ssc install hprescott
hprescott y, smooth(1600)      // HP
// ssc install bking
bking y, minperiod(6) maxperiod(32) stub(bk)
// ssc install cfitzroy
// stl: via Python/R (nao disponivel nativamente no Stata)
\end{lstlisting}

\subsection*{R}

\begin{lstlisting}[language=R, basicstyle=\ttfamily\small,
  frame=single, numbers=none, backgroundcolor=\color{codbg},
  keywordstyle=\color{keyword}\bfseries,
  stringstyle=\color{string},
  commentstyle=\color{comment}\itshape]
library(mFilter)
hp <- hpfilter(y, freq = 1600)
bk <- bkfilter(y, pl = 6, pu = 32, nfix = 12)
cf <- cffilter(y, pl = 6, pu = 32)
library(stats)
stl(ts(y, frequency=12), s.window="periodic")
\end{lstlisting}

\subsection*{Python}

\begin{lstlisting}[language=Python, basicstyle=\ttfamily\small,
  frame=single, numbers=none, backgroundcolor=\color{codbg},
  keywordstyle=\color{keyword}\bfseries,
  stringstyle=\color{string},
  commentstyle=\color{comment}\itshape]
from statsmodels.tsa.filters.hp_filter import hpfilter
from statsmodels.tsa.filters.bk_filter import bkfilter
from statsmodels.tsa.seasonal import STL
_, trend = hpfilter(y, lamb=1600)
cycle_bk = bkfilter(y, low=6, high=32, K=12)
stl = STL(y, period=12).fit()
\end{lstlisting}

% ─────────────────────────────────────────────────────────────────────────────
\section{PSM e Controle Sintético}

\subsection*{Hayashi}

\begin{lstlisting}
let m_psm = psm(y ~ d + x1, df, k=1, caliper=0.05, boot=200)
display m_psm

// Controle Sintético (painel)
let m_sc = synth("y", 1, 11, df, id="unit", time="year")
display m_sc
\end{lstlisting}

\subsection*{Stata}

\begin{lstlisting}[language={}, basicstyle=\ttfamily\small,
  frame=single, numbers=none, backgroundcolor=\color{codbg}]
// ssc install psmatch2
psmatch2 d x1, outcome(y) neighbor(1) caliper(0.05)
// ssc install synth
synth y x1, trunit(1) trperiod(11) xperiod(1/10)
\end{lstlisting}

\subsection*{R}

\begin{lstlisting}[language=R, basicstyle=\ttfamily\small,
  frame=single, numbers=none, backgroundcolor=\color{codbg},
  keywordstyle=\color{keyword}\bfseries,
  stringstyle=\color{string},
  commentstyle=\color{comment}\itshape]
library(MatchIt)
m <- matchit(d ~ x1, data=df, method="nearest",
             caliper=0.05)
library(Synth)
dataprep.out <- dataprep(...)    # setup
synth(dataprep.out)
\end{lstlisting}

\subsection*{Python}

\begin{lstlisting}[language=Python, basicstyle=\ttfamily\small,
  frame=single, numbers=none, backgroundcolor=\color{codbg},
  keywordstyle=\color{keyword}\bfseries,
  stringstyle=\color{string},
  commentstyle=\color{comment}\itshape]
# PSM:
from sklearn.linear_model import LogisticRegression
# calcular propensity scores e fazer matching manual
# Controle Sintético:
# pip install pysynth
from pysynth import Synth
\end{lstlisting}

% ─────────────────────────────────────────────────────────────────────────────
\section{Resumo comparativo}

{\footnotesize\setlength{\tabcolsep}{3pt}%
\begin{longtable}{lllll}
\toprule
\textbf{Tarefa} & \textbf{Hayashi} & \textbf{Stata} & \textbf{R} & \textbf{Python} \\
\midrule
\endfirsthead
\multicolumn{5}{l}{\footnotesize\itshape (continuação)} \\[2pt]
\toprule
\textbf{Tarefa} & \textbf{Hayashi} & \textbf{Stata} & \textbf{R} & \textbf{Python} \\
\midrule
\endhead
\midrule
\multicolumn{5}{r}{\footnotesize\itshape (continua na próxima página)} \\
\endfoot
\bottomrule
\endlastfoot
%
\multicolumn{5}{l}{\textit{Regressão}} \\
OLS             & \texttt{ols()}        & \texttt{reg}          & \texttt{lm()}          & \texttt{smf.ols()} \\
IV/2SLS         & \texttt{iv()}         & \texttt{ivregress}    & \texttt{ivreg()}       & \texttt{IV2SLS()} \\
GMM             & \texttt{gmm()}        & \texttt{gmm}          & \texttt{gmm()}         & \texttt{LinearIVGMM()} \\
WLS             & \texttt{wls()}        & \texttt{vwls}         & \texttt{lm(weights=)}  & \texttt{WLS()} \\
Rolling OLS     & \texttt{rols()}       & \texttt{rolling()}    & \texttt{rollregres()}  & \texttt{RollingOLS()} \\
LASSO           & \texttt{lasso()}      & \texttt{lasso linear} & \texttt{glmnet(a=1)}   & \texttt{Lasso()} \\
Ridge           & \texttt{ridge\_reg()} & —                     & \texttt{glmnet(a=0)}   & \texttt{Ridge()} \\
Elastic Net     & \texttt{enet()}       & —                     & \texttt{glmnet(a=.5)}  & \texttt{ElasticNet()} \\
Bootstrap       & \texttt{bootstrap()}  & \texttt{bootstrap:}   & \texttt{boot()}        & \texttt{resample()} \\
Bootstrap SE    & \texttt{bootse()}     & —                     & —                      & —                     \\
LOWESS          & \texttt{lowess()}     & \texttt{lowess}       & \texttt{lowess()}      & \texttt{lowess()} \\
Quant. reg.     & \texttt{qreg()}       & \texttt{qreg}         & \texttt{rq()}          & \texttt{QuantReg()} \\
SUR             & \texttt{sur()}        & \texttt{sureg}        & \texttt{systemfit()}   & \texttt{SUR()} \\
\midrule
\multicolumn{5}{l}{\textit{Dados em painel}} \\
FE (within)     & \texttt{fe()}         & \texttt{xtreg, fe}    & \texttt{plm(within)}   & \texttt{PanelOLS()} \\
RE (GLS)        & \texttt{re()}         & \texttt{xtreg, re}    & \texttt{plm(random)}   & \texttt{RandomEffects()} \\
Hausman         & \texttt{hausman()}    & \texttt{hausman}      & \texttt{phtest()}      & manual \\
Hausman robusto & \texttt{hausman\_robust()} & manual           & \texttt{phtest(robust)} & manual \\
AB-GMM          & \texttt{ab()}         & \texttt{xtabond2}     & \texttt{pgmm()}        & \texttt{FirstDifferenceGMM()} \\
PCSE            & \texttt{pcse()}       & \texttt{xtpcse}       & \texttt{pcse()}        & — \\
GEE             & \texttt{gee()}        & \texttt{xtgee}        & \texttt{geeglm()}      & \texttt{GEE()} \\
Mundlak (CRE)   & \texttt{mundlak()}    & manual                & \texttt{mundlak()}     & manual \\
xtgls           & \texttt{xtgls()}      & \texttt{xtgls}        & \texttt{pggls()}       & — \\
Modelo misto    & \texttt{mixedlm()}    & \texttt{mixed}        & \texttt{lme4::lmer()}  & \texttt{MixedLM()} \\
Logit FE        & \texttt{clogit()}     & \texttt{clogit}       & \texttt{clogit()}      & \texttt{ConditionalLogit()} \\
Resumo painel   & \texttt{xtsum()}      & \texttt{xtsum}        & —                      & —                     \\
\midrule
\multicolumn{5}{l}{\textit{Variáveis dependentes especiais}} \\
Logit           & \texttt{logit()}      & \texttt{logit}        & \texttt{glm(logit)}    & \texttt{smf.logit()} \\
Probit          & \texttt{probit()}     & \texttt{probit}       & \texttt{glm(probit)}   & \texttt{smf.probit()} \\
Logit ordenado  & \texttt{ologit()}     & \texttt{ologit}       & \texttt{polr()}        & \texttt{OrderedModel} \\
Probit ordenado & \texttt{oprobit()}    & \texttt{oprobit}      & \texttt{polr()}        & \texttt{OrderedModel} \\
Logit multinom. & \texttt{mlogit()}     & \texttt{mlogit}       & \texttt{multinom()}    & \texttt{MNLogit()} \\
Poisson         & \texttt{poisson()}    & \texttt{poisson}      & \texttt{glm(poisson)}  & \texttt{smf.poisson()} \\
Neg. Binomial   & \texttt{nbreg()}      & \texttt{nbreg}        & \texttt{glm.nb()}      & \texttt{smf.negativebinomial()} \\
Zero-inflated   & \texttt{zinb()}       & \texttt{zinb}         & \texttt{zeroinfl()}    & \texttt{ZeroInflatedNB} \\
Tobit           & \texttt{tobit()}      & \texttt{tobit}        & \texttt{AER::tobit()}  & manual \\
Heckman         & \texttt{heckman()}    & \texttt{heckman}      & \texttt{heckit()}      & manual \\
GLM geral       & \texttt{glm()}        & \texttt{glm}          & \texttt{glm()}         & \texttt{GLM()} \\
Ef. marginais   & \texttt{margins()}    & \texttt{margins}      & \texttt{margins()}     & \texttt{get\_margeff()} \\
\midrule
\multicolumn{5}{l}{\textit{Inferência causal}} \\
DiD             & \texttt{did()}        & \texttt{diff}         & \texttt{lm()}          & \texttt{smf.ols()} \\
RDD             & \texttt{rd()}         & \texttt{rdrobust}     & \texttt{rdrobust()}    & \texttt{rdd()} \\
Fuzzy RDD       & \texttt{fuzzy\_rd()}   & \texttt{rdrobust}     & \texttt{rdrobust()}    & \texttt{rdd()} \\
PSM             & \texttt{psm()}        & \texttt{psmatch2}     & \texttt{matchit()}     & manual \\
Controle sint.  & \texttt{synth()}      & \texttt{synth}        & \texttt{Synth}         & \texttt{pysynth} \\
\midrule
\multicolumn{5}{l}{\textit{Sobrevivência}} \\
Kaplan-Meier    & \texttt{km()}         & \texttt{sts}          & \texttt{survfit()}     & \texttt{KaplanMeierFitter} \\
Cox PH          & \texttt{cox()}        & \texttt{stcox}        & \texttt{coxph()}       & \texttt{CoxPHFitter} \\
\midrule
\multicolumn{5}{l}{\textit{Séries temporais}} \\
AR/MA/ARMA      & \texttt{arima()}      & \texttt{arima}        & \texttt{arima()}       & \texttt{ARIMA()} \\
AR($p$) OLS     & \texttt{autoreg()}    & \texttt{arima}        & \texttt{arima()}       & \texttt{AutoReg()} \\
ARDL($p$,$q$)   & \texttt{ardl()}       & \texttt{ardl}         & \texttt{dynlm()}       & \texttt{ardl()} \\
GARCH           & \texttt{garch()}      & \texttt{arch}         & \texttt{garch()}       & \texttt{arch()} \\
Raiz unit. (ADF) & \texttt{adf()}       & \texttt{dfuller}      & \texttt{adf.test()}    & \texttt{adfuller()} \\
Phillips-Perron & \texttt{phillips\_perron()} & \texttt{pperron} & \texttt{PP.test()}   & \texttt{PhillipsPerron()} \\
Zivot-Andrews   & \texttt{zivot()}      & \texttt{zandrews}     & \texttt{ur.za()}       & \texttt{ZivotAndrews()} \\
Cointegração    & \texttt{johansen()}   & \texttt{vecrank}      & \texttt{ca.jo()}       & \texttt{coint\_johansen} \\
VAR             & \texttt{var()}        & \texttt{var}          & \texttt{VAR()}         & \texttt{VAR()} \\
IRF (reduzida)  & \texttt{irf()}        & \texttt{irf create}   & \texttt{irf()}         & \texttt{irf().plot()} \\
SVAR Cholesky   & \texttt{svar(chol.)}  & \texttt{svar}         & \texttt{id.chol()}     & manual \\
SVAR BQ         & \texttt{svar(long.)}  & \texttt{svar}         & \texttt{BQ()}          & manual \\
IRF estrutural  & \texttt{sirf()}       & \texttt{irf create}   & \texttt{irf()}         & manual \\
Markov Switch.  & \texttt{markov()}     & \texttt{mswitch}      & \texttt{msmFit()}      & \texttt{MarkovRegression} \\
\midrule
\multicolumn{5}{l}{\textit{Filtros e decomposição}} \\
HP              & \texttt{hprescott()}  & \texttt{hprescott}$^*$ & \texttt{hpfilter()}  & \texttt{hpfilter()} \\
Baxter-King     & \texttt{baxter\_king()} & \texttt{bking}$^*$  & \texttt{bkfilter()}   & \texttt{bkfilter()} \\
Holt-Winters    & \texttt{holtwinters()} & \texttt{tssmooth}   & \texttt{HoltWinters()} & \texttt{ExponentialSmoothing()} \\
STL             & \texttt{stl()}        & —                     & \texttt{stl()}         & \texttt{STL()} \\
CF filter       & \texttt{christiano\_fitzgerald()} & \texttt{cfitzroy} & \texttt{cffilter()} & — \\
\midrule
\multicolumn{5}{l}{\textit{Análise multivariada}} \\
PCA             & \texttt{pca()}        & \texttt{pca}          & \texttt{prcomp()}      & \texttt{PCA()} \\
Análise fatorial & \texttt{factor()}    & \texttt{factor}       & \texttt{factanal()}    & \texttt{FactorAnalysis()} \\
Cor. canônica   & \texttt{cancorr()}    & \texttt{canon}        & \texttt{cancor()}      & \texttt{CCA()} \\
MANOVA          & \texttt{manova()}     & \texttt{manova}       & \texttt{manova()}      & \texttt{MANOVA()} \\
\midrule
\multicolumn{5}{l}{\textit{Agrupamento e redução de dimensionalidade}} \\
K-Means         & \texttt{kmeans()}     & \texttt{cluster kmeans} & \texttt{kmeans()}    & \texttt{KMeans()} \\
Hierárquico     & \texttt{hclust()}     & \texttt{cluster single} & \texttt{hclust()}    & \texttt{AgglomerativeClustering()} \\
DBSCAN          & \texttt{dbscan()}     & —                     & \texttt{dbscan()}      & \texttt{DBSCAN()} \\
t-SNE           & \texttt{tsne()}       & —                     & \texttt{tsne()}        & \texttt{TSNE()} \\
UMAP            & \texttt{umap()}       & —                     & \texttt{umap()}        & \texttt{UMAP()} \\
Biplot PCA      & \texttt{biplot()}     & \texttt{pca}          & \texttt{biplot()}      & manual \\
\midrule
\multicolumn{5}{l}{\textit{Pós-estimação}} \\
Coefs. tidy     & \texttt{tidy()}       & —                     & \texttt{broom::tidy()} & \texttt{statsmodels summary} \\
Glance          & \texttt{glance()}     & —                     & \texttt{broom::glance()} & manual \\
\midrule
\multicolumn{5}{l}{\textit{Inferência e testes}} \\
Teste $t$       & \texttt{ttest()}      & \texttt{ttest}        & \texttt{t.test()}      & \texttt{ttest\_ind()} \\
ANOVA           & \texttt{anova()}      & \texttt{anova}        & \texttt{aov()}         & \texttt{f\_oneway()} \\
Kruskal-Wallis  & \texttt{kruskal\_wallis()} & \texttt{kwallis} & \texttt{kruskal.test()} & \texttt{kruskal()} \\
Normalidade     & \texttt{swilk()}      & \texttt{swilk}        & \texttt{shapiro.test()} & \texttt{shapiro()} \\
Teste conjunto  & \texttt{testparm()}   & \texttt{testparm}     & \texttt{linearHypothesis()} & \texttt{wald\_test()} \\
Tabela          & \texttt{esttab()}     & \texttt{esttab}$^\dagger$ & \texttt{stargazer()} & \texttt{summary\_col()} \\
Armazenar modelo& \texttt{eststo()}     & \texttt{eststo}       & —                      & —                     \\
Limpar modelos  & \texttt{estclear()}   & —                     & —                      & —                     \\
\end{longtable}}

\noindent
$^*$ HP e BK no Stata requerem \texttt{ssc install hprescott} e \texttt{ssc install bking}.\\
$^\dagger$ \texttt{esttab} no Stata requer \texttt{ssc install estout}.\\[4pt]
\noindent\textbf{Nota sobre geradores aleatórios:} \texttt{rnormal()} corresponde ao
gerador normal padrão usado nos exemplos de Stata, R e Python. Use
\texttt{runiform()} quando o DGP exigir um sorteio $U(0,1)$.

% ─────────────────────────────────────────────────────────────────────────────
\section{Matriz de validação}
\label{sec:matriz_validacao}

Além dos exemplos de código deste apêndice, todo \hayashi{} é validado
sistematicamente contra implementações de referência. O programa de
validação está no diretório \texttt{validation/} e é executado com:

\begin{lstlisting}
hay validate
\end{lstlisting}

A Tabela~\ref{tab:matriz_validacao} lista os casos de validação atuais, com a
família de estimadores, a fonte dos dados e as referências (R, Python e/ou
Stata) usadas na comparação. Os detalhes completos --- tolerâncias, notas e
status --- estão em \texttt{validation/MATRIX.md}.

{\footnotesize\setlength{\tabcolsep}{4pt}%
\begin{longtable}{llll}
\toprule
\label{tab:matriz_validacao}
\textbf{ID do Caso} & \textbf{Família} & \textbf{Dataset/Fonte} & \textbf{Referências (R/Python/Stata)} \\
\midrule
\endfirsthead
\multicolumn{4}{l}{\footnotesize\itshape (continuação)} \\[2pt]
\toprule
\textbf{ID do Caso} & \textbf{Família} & \textbf{Dataset/Fonte} & \textbf{Referências (R/Python/Stata)} \\
\midrule
\endhead
\midrule
\multicolumn{4}{r}{\footnotesize\itshape (continua na próxima página)} \\
\endfoot
\bottomrule
\endlastfoot
%
\texttt{ab\_grunfeld} & \texttt{ab} & \texttt{wooldridge::grunfeld} & R \\
\texttt{ar\_book} & \texttt{arima} & \texttt{simulated\_ar1} & R, Python \\
\texttt{ardl\_gdp} & \texttt{ardl} & \texttt{statsmodels::macrodata} & R, Python \\
\texttt{arima\_book} & \texttt{arima} & \texttt{simulated\_rw} & R, Python \\
\texttt{arima\_gdp} & \texttt{arima} & \texttt{statsmodels::macrodata} & R, Python \\
\texttt{arma\_book} & \texttt{arima} & \texttt{simulated\_arma11} & R, Python \\
\texttt{autoreg\_gdp} & \texttt{autoreg} & \texttt{statsmodels::macrodata} & R, Python \\
\texttt{coint\_book} & \texttt{vecm} & \texttt{simulated\_cointegrated} & R, Python \\
\texttt{cox\_heart} & \texttt{cox} & \texttt{statsmodels::heart} & R, Python \\
\texttt{did\_kielmc} & \texttt{did} & \texttt{wooldridge::kielmc} & R, Python \\
\texttt{ets\_gdp} & \texttt{ets} & \texttt{statsmodels::macrodata} & R, Python \\
\texttt{garch\_book} & \texttt{garch} & \texttt{simulated\_garch11} & Python \\
\texttt{garch\_nyse} & \texttt{garch} & \texttt{wooldridge::nyse} & R, Python \\
\texttt{glsar\_hprice1} & \texttt{glsar} & \texttt{wooldridge::hprice1} & R, Python \\
\texttt{gmm\_card} & \texttt{gmm} & \texttt{wooldridge::card} & R, Python \\
\texttt{heckman\_mroz} & \texttt{heckman} & \texttt{wooldridge::mroz} & R, Python \\
\texttt{iv\_card} & \texttt{iv} & \texttt{wooldridge::card} & R, Python \\
\texttt{lasso\_hprice1} & \texttt{lasso} & \texttt{wooldridge::hprice1} & R, Python \\
\texttt{logit\_mroz} & \texttt{logit} & \texttt{wooldridge::mroz} & R, Python \\
\texttt{ma\_book} & \texttt{arima} & \texttt{simulated\_ma1} & R, Python \\
\texttt{ologit\_beauty} & \texttt{logit} & \texttt{wooldridge::beauty} & R, Python \\
\texttt{ols\_wooldridge\_wage1} & \texttt{ols} & \texttt{wooldridge::wage1} & R, Python \\
\texttt{panel\_fe\_grunfeld} & \texttt{panel\_fe} & \texttt{wooldridge::grunfeld} & R, Python \\
\texttt{poisson\_fertil2} & \texttt{poisson} & \texttt{wooldridge::fertil2} & R, Python \\
\texttt{probit\_mroz} & \texttt{probit} & \texttt{wooldridge::mroz} & R, Python \\
\texttt{psm\_jtrain3} & \texttt{psm} & \texttt{wooldridge::jtrain3} & R, Python \\
\texttt{quantile\_wage1} & \texttt{qreg} & \texttt{wooldridge::wage1} & R, Python \\
\texttt{re\_grunfeld} & \texttt{re} & \texttt{grunfeld} & R, Python \\
\texttt{rdd\_book} & \texttt{rdd} & \texttt{rdd\_book} & R, Python \\
\texttt{synth\_smoking} & \texttt{synth} & \texttt{synth\_smoking} & R, Python \\
\texttt{tobit\_mroz} & \texttt{tobit} & \texttt{wooldridge::mroz} & R, Python \\
\texttt{var\_book} & \texttt{var} & \texttt{simulated\_var1} & R, Python \\
\texttt{var\_macro} & \texttt{var} & \texttt{statsmodels::macrodata} & R, Python \\
\texttt{wls\_hprice1} & \texttt{wls} & \texttt{wooldridge::hprice1} & R, Python \\
\end{longtable}}
